\documentclass[aspectratio=169,11pt]{beamer}

% =========================================================
% Packages
% =========================================================
\usepackage{amsmath,amssymb,mathtools,bm}
\usepackage{booktabs}
\usepackage{array}
\usepackage{ragged2e}
\usepackage{microtype}
\usepackage{graphicx}
\usepackage{xcolor}

% =========================================================
% Theme
% =========================================================
\usetheme{Boadilla}
\usecolortheme{default}
\setbeamertemplate{navigation symbols}{}
\setbeamertemplate{footline}[frame number]

\definecolor{DeepBlue}{RGB}{28,66,120}
\definecolor{DarkGray}{RGB}{55,55,55}
\definecolor{SoftGray}{RGB}{242,242,242}
\definecolor{DarkRed}{RGB}{155,45,45}
\definecolor{DarkGreen}{RGB}{35,110,75}

\setbeamercolor{structure}{fg=DeepBlue}
\setbeamercolor{frametitle}{fg=DeepBlue}
\setbeamercolor{title}{fg=DeepBlue}
\setbeamercolor{block title}{bg=DeepBlue!12,fg=DeepBlue}
\setbeamercolor{block body}{bg=DeepBlue!3,fg=black}
\setbeamercolor{alerted text}{fg=DarkRed}

% =========================================================
% Layout tuning (compact vertical rhythm)
% =========================================================
\setbeamersize{text margin left=1.1em,text margin right=1.1em}

% Many slides contain several display equations; the default
% display skips are the single largest source of overfull frames.
\AtBeginDocument{%
  \setlength{\abovedisplayskip}{3.5pt plus 2pt minus 1.5pt}%
  \setlength{\belowdisplayskip}{3.5pt plus 2pt minus 1.5pt}%
  \setlength{\abovedisplayshortskip}{1.5pt plus 1pt minus 1pt}%
  \setlength{\belowdisplayshortskip}{2.5pt plus 1pt minus 1pt}%
  \setlength{\jot}{2pt}%
}

% Tighter blocks: trim the space around the box, not its inner padding
\addtobeamertemplate{block begin}{\vspace{-0.45em}}{}
\addtobeamertemplate{block end}{}{\vspace{-0.35em}}
\addtobeamertemplate{block alerted begin}{\vspace{-0.45em}}{}
\addtobeamertemplate{block alerted end}{}{\vspace{-0.35em}}

% Tighter lists
\setlength{\leftmargini}{1.5em}
\setbeamertemplate{itemize/enumerate body begin}{\vspace{-0.2em}}
\setbeamertemplate{itemize/enumerate body end}{\vspace{-0.2em}}

% Slightly tighter frame titles
\setbeamerfont{frametitle}{size=\large,series=\bfseries}

% Ragged-right table columns; narrow columns look much calmer when
% they are not hyphenated on nearly every line.
\newcommand{\NoHyph}{\hyphenpenalty=10000\exhyphenpenalty=10000\relax}
\newcolumntype{L}[1]{>{\RaggedRight\NoHyph\arraybackslash}p{#1}}
\newcolumntype{B}[1]{>{\RaggedRight\NoHyph\bfseries\arraybackslash}p{#1}}

% =========================================================
% Commands
% =========================================================
\newcommand{\E}{\mathbb{E}}
\newcommand{\Pp}{\mathbb{P}}
\newcommand{\ind}{\mathbf{1}}
\newcommand{\ATE}{\mathrm{ATE}}
\newcommand{\ATT}{\mathrm{ATT}}

\newcommand{\Question}{\textcolor{DarkRed}{\textbf{Question}}}
\newcommand{\Spec}{\textcolor{DeepBlue}{\textbf{Specification \& Estimand}}}
\newcommand{\Ident}{\textcolor{DarkGreen}{\textbf{Identification}}}
\newcommand{\Est}{\textcolor{violet}{\textbf{Estimation}}}
\newcommand{\Infer}{\textcolor{orange!80!black}{\textbf{Inference}}}

% =========================================================
% Title
% =========================================================
\title{Mostly Harmless Econometrics}
\subtitle{
\textbf{A Course Correction in the History of Econometrics}
\\[0.55em]
Question $\rightarrow$ Specification \& Estimand
$\rightarrow$ Identification
$\rightarrow$ Estimation
$\rightarrow$ Inference
}
\author{Kaiwen Luo}
\institute{Washington University in St.Louis}
\date{\today}

\begin{document}

% =========================================================
\begin{frame}
    \titlepage
\end{frame}

% =========================================================
\section{Why Econometrics?}

\begin{frame}{A familiar way to learn econometrics}

Many traditional econometrics courses are organized around estimators:
\[
\text{OLS}
\rightarrow
\text{GLS}
\rightarrow
\text{MLE}
\rightarrow
\text{IV/2SLS}
\rightarrow
\text{Panel FE}
\rightarrow
\text{Discrete Choice}.
\]

\pause

\vspace{0.3em}

This training is valuable. It teaches us how to solve problems such as

\[
\min_{\beta}\sum_i (Y_i-X_i'\beta)^2,
\]
how to derive

\[
\sqrt{n}(\hat\beta-\beta_0)
\overset{d}{\longrightarrow}
N(0,V),
\]
and how to construct standard errors.

\pause

\vspace{0.35em}

\begin{alertblock}{But there is a more fundamental question}
Before asking how to estimate $\beta$, we should ask:
\[
\boxed{\text{Why is } \beta \text{ the object we want in the first place?}}
\]
\end{alertblock}

\end{frame}

% =========================================================
\begin{frame}{The econometric workflow}
\begin{center}
\large
\Question

\vspace{0.15em}
$\Downarrow$

\Spec

\vspace{0.15em}
$\Downarrow$

\Ident

\vspace{0.15em}
$\Downarrow$

\Est

\vspace{0.15em}
$\Downarrow$

\Infer
\end{center}

\vspace{0.2em}

\begin{block}{The central principle of this course}
An estimator is not a research design.

A regression equation is not automatically a causal model.

A model that can be estimated is not necessarily a model that answers
the question we care about.
\end{block}
\end{frame}

% =========================================================
\begin{frame}{What does each step ask?}

\small

\begin{tabular}{B{0.19\textwidth} L{0.72\textwidth}}
\toprule
Question
&
What economic or causal question are we trying to answer?
\\[0.5em]

Specification
&
What are the units, outcomes, choices, treatments, states,
constraints, and behavioral relationships?
\\[0.5em]

Estimand
&
What precise population object represents the answer to our question?
\\[0.5em]

Identification
&
Under what assumptions is that object uniquely determined by the
population distribution of observed data?
\\[0.5em]

Estimation
&
Given a finite sample, how do we construct an estimator of that object?
\\[0.5em]

Inference
&
How uncertain is our estimate because we observe only a finite sample?
\\
\bottomrule
\end{tabular}

\vspace{0.4em}

\begin{alertblock}{Do not reverse the order}
\[
\boxed{
\text{``I can estimate this regression''}
\;\not\Rightarrow\;
\text{``this regression answers my question.''}
}
\]
\end{alertblock}

\end{frame}

% =========================================================
\begin{frame}{Specification and identification are different}

Suppose a model is indexed by a parameter $\theta$ and implies an
observable-data distribution $P_\theta$.

A classical definition of identification is:
\[
P_{\theta_1}=P_{\theta_2}
\quad\Longrightarrow\quad
\theta_1=\theta_2.
\]

\vspace{0.3em}

\pause

\textbf{Specification} comes first:
\[
\boxed{
\text{Which family } \{P_\theta:\theta\in\Theta\}
\text{ are we willing to consider?}
}
\]

\begin{center}
\textcolor{DeepBlue}{\textit{``All models are wrong, but some are useful.''}}
\quad
\textcolor{DarkGray}{--- George E.\,P. Box}
\end{center}

\pause

Then \textbf{identification} asks:
\[
\boxed{
\text{Within that model, can different values of }\theta
\text{ generate the same observable data?}
}
\]

\vspace{0.35em}

\begin{alertblock}{Important}
Identification can be perfect \emph{within a badly chosen model}.

That does not rescue the specification.
\end{alertblock}

\end{frame}

% =========================================================
\begin{frame}{Two examples for today}

We will study the same workflow through two very different traditions.

\vspace{0.5em}

\begin{columns}[T]

\column{0.48\textwidth}

\begin{block}{Example 1: Treatment effects}
\textbf{Question:}
Does vaccination reduce COVID infection?

\vspace{0.25em}

Language:
\[
Y(1),Y(0),\ATE
\]

Main flavor:

\begin{center}
\textbf{design-based / reduced-form}
\end{center}

\end{block}

\column{0.48\textwidth}

\begin{block}{Example 2: Binary choice}
\textbf{Question:}
How will consumers respond to a price change?

\vspace{0.25em}

Language:
\[
U_{ij},\theta,\varepsilon_{ij}
\]

Main flavor:

\begin{center}
\textbf{structural}
\end{center}

\end{block}

\end{columns}

\vspace{0.4em}

The methods look very different.

\pause

But the discipline of reasoning should be the same.

\end{frame}

% =========================================================
\section{Example I: ATE and Vaccination}

\begin{frame}{Example I: Start from the question}

Suppose someone asks:

\begin{center}
\Large
\textbf{``What is the effect of vaccination on COVID infection?''}
\end{center}

\vspace{0.5em}

This sounds like a well-defined causal question.

\pause

It is not.

\vspace{0.35em}

We still need to specify:
\[
\boxed{
\begin{array}{c}
\text{Who?}\\
\text{Which vaccine?}\\
\text{Vaccinated when?}\\
\text{Outcome measured when?}\\
\text{Infection or severe disease?}\\
\text{One person's vaccination or a vaccination campaign?}
\end{array}}
\]

\vspace{0.25em}

\begin{alertblock}{Lesson}
Natural-language questions are usually not yet econometric estimands.
\end{alertblock}

\end{frame}

% =========================================================
\begin{frame}{A first specification: individual treatment}
Let $D_i\in\{0,1\}$ indicate whether individual $i$ is vaccinated, and define
the potential outcomes
\[
Y_i(1) = \text{infection outcome if }i\text{ is vaccinated},
\]
\[
Y_i(0) = \text{infection outcome if }i\text{ is not vaccinated}.
\]
Then an apparently natural estimand is
\[
\boxed{\ATE = \E[Y_i(1)-Y_i(0)].}
\]
Now we have translated a verbal question into a mathematical object.
\pause

\begin{alertblock}{But have we translated the question correctly?}
That is a \textbf{specification question}, not an estimation question.
\end{alertblock}
\end{frame}

% =========================================================
\begin{frame}{What is hidden inside $Y_i(1)$ and $Y_i(0)$?}

Writing $Y_i(d)$ implicitly assumes that once we specify $i$'s
treatment $d$, the potential outcome is well defined.

A standard version of SUTVA requires, among other things:

\vspace{0.3em}

\begin{enumerate}
    \item no relevant hidden versions of treatment;
    \item no interference across units.
\end{enumerate}

The second condition means:
\[
Y_i
\text{ depends on } D_i,
\quad
\text{but not on } D_j
\text{ for }j\neq i.
\]

Equivalently,
\[
Y_i(\bm D)
=
Y_i(D_i).
\]

\pause

\vspace{0.25em}

\begin{alertblock}{This is not merely technical notation}
It is an economic and scientific statement about how the world works.
\end{alertblock}

\end{frame}

% =========================================================
\begin{frame}{Vaccination exposes the specification problem}
For an infectious disease, individual $i$'s infection risk can depend on $D_i$,
and also on the vaccination choices of the people around $i$. For example,
\[
Y_i = Y_i(D_i,D_{-i}).
\]
Or, using a local vaccination rate $G_i=\frac{1}{|N_i|}\sum_{j\in N_i}D_j$,
we may instead need
\[
\boxed{Y_i(d,g).}
\]
If vaccination changes transmission, then
\[
Y_i(0,g=0.8) \neq Y_i(0,g=0.2)
\]
may be perfectly possible.

\begin{alertblock}{Therefore}
The simple pair $\{Y_i(1),Y_i(0)\}$ may be too small a model for the
substantive question.
\end{alertblock}
\end{frame}

% =========================================================
\begin{frame}{The research question determines the estimand}
Once interference is allowed, several distinct questions appear.

\textbf{Direct effect at exposure level $g$:}
\[
DE(g) = \E[Y_i(1,g)-Y_i(0,g)].
\]
\textbf{Spillover effect for an untreated person:}
\[
SE(g_1,g_0) = \E[Y_i(0,g_1)-Y_i(0,g_0)].
\]
\textbf{Policy effect:}
\[
PE(p_1,p_0) = \E[Y_i(\bm D(p_1))-Y_i(\bm D(p_0))],
\]
where $p$ indexes a population vaccination policy.

\begin{alertblock}{These are different causal questions}
There is no unique object called ``the effect of vaccination.''
\end{alertblock}
\end{frame}

% =========================================================
\begin{frame}{A perfectly executed experiment can answer the wrong question}

Suppose we randomly vaccinate half of the individuals:
\[
D_i\perp \{Y_i(1),Y_i(0)\}.
\]

Under the simple no-interference model,
\[
\E[Y_i\mid D_i=1]
-
\E[Y_i\mid D_i=0]
=
\ATE.
\]

Excellent identification.

\vspace{0.3em}

\pause

But suppose the policy question is:

\begin{center}
\textit{What happens when population coverage rises from 20\% to 80\%?}
\end{center}

Then an individually randomized experiment conducted around 50\% coverage
need not identify that policy effect.

\vspace{0.4em}

\begin{alertblock}{The deeper lesson}
\[
\boxed{
\text{Strong identification}
+
\text{wrong estimand}
=
\text{a precise answer to the wrong question}.
}
\]
\end{alertblock}

\end{frame}

% =========================================================
\begin{frame}{Only now do we discuss identification}

Return temporarily to the simple potential-outcome model:
\[
Y_i(1),\qquad Y_i(0).
\]

Observed outcome:
\[
Y_i
=
D_iY_i(1)+(1-D_i)Y_i(0).
\]

We want

\[
\ATE=\E[Y_i(1)-Y_i(0)].
\]

But for every individual we observe only one potential outcome.

This is the fundamental missing-data problem:
\[
\begin{array}{c|cc}
 & Y_i(1) & Y_i(0)\\
\hline
D_i=1 & \checkmark & ?\\
D_i=0 & ? & \checkmark
\end{array}
\]

Identification asks:
\[
\boxed{
\text{What assumption allows observed outcomes to reconstruct
the missing counterfactual?}
}
\]

\end{frame}

% =========================================================
\begin{frame}{Random assignment is an identification strategy}
Suppose
\[
D_i \perp \left(Y_i(1),Y_i(0)\right).
\]
Then
\[
\E[Y_i(1)] = \E[Y_i\mid D_i=1],
\qquad
\E[Y_i(0)] = \E[Y_i\mid D_i=0],
\]
and therefore
\[
\boxed{\ATE = \E[Y_i\mid D_i=1]-\E[Y_i\mid D_i=0].}
\]
Notice what happened: the causal estimand
\[
\underbrace{\E[Y_i(1)-Y_i(0)]}_{\text{causal estimand}}
\]
became equal, under an identification assumption, to an observable quantity
\[
\underbrace{\E[Y_i\mid D_i=1]-\E[Y_i\mid D_i=0]}_{\text{observable population quantity}}.
\]
\end{frame}

% =========================================================
\begin{frame}{Identification is not estimation}
Once identification is established,
\[
\ATE = \E[Y\mid D=1]-\E[Y\mid D=0],
\]
we can estimate it using
\[
\boxed{\widehat{\ATE} = \bar Y_1-\bar Y_0.}
\]
This is \textbf{estimation}. We can then derive
\[
\mathrm{Var}(\widehat{\ATE}) = \frac{\sigma_1^2}{n_1}+\frac{\sigma_0^2}{n_0},
\]
and construct a confidence interval. This is \textbf{inference}.

\begin{alertblock}{Three logically different statements}
\[
\begin{array}{ll}
\text{Identification:}
&
\ATE = \E[Y\mid D=1]-\E[Y\mid D=0].
\\[0.25em]
\text{Estimation:}
&
\widehat{\ATE}=\bar Y_1-\bar Y_0.
\\[0.25em]
\text{Inference:}
&
\widehat{\ATE}\pm 1.96\,SE(\widehat{\ATE}).
\end{array}
\]
\end{alertblock}
\end{frame}

% =========================================================
\begin{frame}{Standard errors do not measure identification uncertainty}

Suppose the estimate is

\[
\widehat{\ATE}=-0.08,
\qquad
SE(\widehat{\ATE})=0.01.
\]

This tells us that sampling noise is relatively small.

It does \textbf{not} tell us whether

\[
D_i\perp(Y_i(1),Y_i(0))
\]
is credible.

It does not tell us whether SUTVA holds.

It does not tell us whether we selected the correct outcome horizon.

It does not tell us whether the estimand corresponds to the policy question.

\vspace{0.35em}

\begin{block}{Inference is conditional}
Conventional confidence intervals quantify sampling uncertainty
\emph{conditional on the maintained model and identification assumptions}.
\end{block}

\end{frame}

% =========================================================
\begin{frame}{Example I: the complete workflow}

\small

\begin{tabular}{B{0.22\textwidth} L{0.68\textwidth}}
\toprule
Question
&
What aspect of vaccination are we trying to evaluate?
\\[0.5em]

Specification
&
Individual treatment?
Population coverage?
Network exposure?
What outcome and horizon?
\\[0.5em]

Estimand
&
ATE?
Direct effect?
Spillover effect?
Policy effect?
\\[0.5em]

Identification
&
Random assignment?
No interference?
Cluster randomization?
Two-stage saturation design?
\\[0.5em]

Estimation
&
Difference in means?
Regression adjustment?
IPW?
\\[0.5em]

Inference
&
Individual-level SE?
Clustered SE?
Randomization inference?
\\
\bottomrule
\end{tabular}

\vspace{0.4em}

\begin{center}
\Large
\boxed{\text{The estimator is near the end of the argument.}}
\end{center}

\end{frame}

% =========================================================
\section{Example II: Binary Choice}

\begin{frame}{Example II: A structural question}

Now consider a different research question.

Suppose consumers must choose between products $A$ and $B$.

We want to know:

\begin{center}
\Large
\textbf{How would demand change if the price of $A$ fell by \$5?}
\end{center}

\vspace{0.4em}

Perhaps we also want:
\[
\text{consumer surplus},
\]

\[
\text{substitution patterns},
\]
or responses to a price never observed in the data.

\vspace{0.35em}

These questions naturally invite a model of

\[
\boxed{\text{how choices are generated}.}
\]

\end{frame}

% =========================================================
\begin{frame}{A standard binary-choice specification}
Suppose utility from the two alternatives is
\[
U_{iA} = X_{iA}'\beta-\alpha p_A+\varepsilon_{iA},
\qquad
U_{iB} = X_{iB}'\beta-\alpha p_B+\varepsilon_{iB}.
\]
Consumer $i$ chooses $A$ if $U_{iA}>U_{iB}$. Define the latent index
\[
\Delta U_i = U_{iA}-U_{iB}.
\]
Then the observed choice is
\[
D_i = \ind\{\Delta U_i>0\}.
\]

\vspace{0.3em}

This is more than a statistical regression.

It specifies a \textbf{decision rule} connecting primitives to behavior.
\end{frame}

% =========================================================
\begin{frame}{From economic structure to an observable probability}

Write

\[
\Delta U_i
=
\Delta X_i'\beta
-
\alpha\Delta p_i
+
\nu_i.
\]

Then

\[
\Pp(D_i=1\mid X_i,p_i)
=
\Pp
\left(
\nu_i
>
-\Delta X_i'\beta+\alpha\Delta p_i
\right).
\]

If $\nu_i$ has the logistic distribution,
\[
\boxed{
\Pp(D_i=1\mid X_i,p_i)
=
\Lambda
\left(
\Delta X_i'\beta-\alpha\Delta p_i
\right).
}
\]

If $\nu_i$ is normal, we obtain a probit model.

\vspace{0.4em}

\begin{alertblock}{The distributional assumption has economic content}
It determines how unobserved utility heterogeneity maps into observed
choice probabilities.
\end{alertblock}

\end{frame}

% =========================================================
\begin{frame}{A hidden assumption: almost everyone has a strict ranking}
Under a standard continuous random-utility specification,
\[
\Pp(U_{iA}=U_{iB})=0.
\]
Almost every consumer strictly prefers one alternative. Thus $D_i=1$ is
interpreted as evidence that $U_{iA}>U_{iB}$.

\vspace{0.3em}

But suppose the empirical setting is different:
\begin{center}
\textbf{A substantial fraction of consumers genuinely regard
$A$ and $B$ as equally good.}
\end{center}

Perhaps the products are almost identical.

Perhaps many consumers explicitly report ``no preference.''

Perhaps the observed purchase is determined by a default,
position on the screen, or random tie-breaking.
\end{frame}

% =========================================================
\begin{frame}{The model can fit the data and still misrepresent the question}
A standard logit can easily generate $\Pp(D_i=A)\approx 0.5$. But it interprets
this as arising from variation in $\nu_i$, and therefore from consumers with
different \emph{strict} latent preferences.

That is conceptually different from saying:
\[
\boxed{\Pp(U_{iA}=U_{iB})>0.}
\]
These models may fit current market shares similarly. But they can imply very
different predictions for
\[
\text{small price changes},
\qquad
\text{defaults},
\qquad
\text{consumer welfare}.
\]

\begin{alertblock}{Specification again comes before estimation}
MLE cannot tell us whether the behavioral model asks the right question.
\end{alertblock}
\end{frame}

% =========================================================
\begin{frame}{An alternative specification with an indifference region}

Suppose latent relative utility is

\[
V_i
=
\Delta X_i'\beta
-
\alpha\Delta p_i
+
\nu_i.
\]

Let $\kappa\geq0$ denote an indifference threshold.

Define the latent preference state

\[
S_i=
\begin{cases}
A, & V_i>\kappa,\\[0.3em]
I, & |V_i|\leq\kappa,\\[0.3em]
B, & V_i<-\kappa.
\end{cases}
\]

Now there are three economically distinct states:
\[
\boxed{
\text{prefer }A,
\qquad
\text{indifferent},
\qquad
\text{prefer }B.
}
\]

If the consumer is forced to make a binary purchase, we still need a
\textbf{tie-breaking rule} for state $I$.

\end{frame}

% =========================================================
\begin{frame}{Observed binary choice now contains two mechanisms}

Suppose an indifferent consumer chooses $A$ with probability $r_i$.

Then

\[
\Pp(D_i=A\mid X_i)
=
\Pp(V_i>\kappa\mid X_i)
+
r_i
\Pp(|V_i|\leq\kappa\mid X_i).
\]

Observed demand therefore reflects:
\[
\boxed{
\text{preference}
+
\text{indifference}
+
\text{tie-breaking}.
}
\]

\vspace{0.35em}

If the true world has all three mechanisms but we estimate a standard
binary logit, we may attribute all observed behavior to preference.

\vspace{0.35em}

That changes the interpretation of $\alpha$ and $\beta$,
and potentially the predicted response to policy.

\end{frame}

% =========================================================
\begin{frame}{Identification now becomes a serious question}
Suppose all we observe for each consumer is $(D_i,X_i,p_i)$.
Can we separately identify
\[
\beta,\qquad \alpha,\qquad \kappa,\qquad r?
\]
Not necessarily. Different combinations of
\[
\text{preference heterogeneity},
\quad
\text{indifference},
\quad
\text{tie-breaking}
\]
may generate the same observed binary choice probabilities.

Additional information may be useful:
\[
\begin{array}{ll}
\text{randomized prices},
&
\text{repeated choices},
\\[0.2em]
\text{randomized defaults},
&
\text{an explicit ``indifferent'' option},
\\[0.2em]
\multicolumn{2}{c}{\text{additional moments or exclusion restrictions}.}
\end{array}
\]
\end{frame}

% =========================================================
\begin{frame}{Identification is not ``the optimizer converged''}
Suppose we write a likelihood
\[
L(\theta) = \prod_{i=1}^{n} P_\theta(D_i\mid X_i),
\]
and obtain $\hat\theta=\arg\max_{\theta}L(\theta)$. The computer reports:
\begin{center}
\texttt{Convergence achieved.}
\end{center}
\pause
That tells us nothing by itself about identification.
Identification is a population property:
\[
P_{\theta_1}(D\mid X) = P_{\theta_2}(D\mid X)
\quad\Longrightarrow\quad
\theta_1=\theta_2?
\]

\begin{alertblock}{Optimization and identification are different problems}
A unique numerical optimum in one sample does not prove that the
economic primitives are identified.
\end{alertblock}
\end{frame}

% =========================================================
\begin{frame}{Why estimate a structural binary-choice model?}
If the primitives are credibly specified and identified, we can ask
counterfactual questions. For example, under a new price $p_A'=p_A-5$,
we can construct
\[
\Pp_{\hat\theta}\left(D_i(p_A')=A\right).
\]
More generally, a structural model may allow us to study
\[
\boxed{\text{policies not directly observed in the historical data}.}
\]
Potential outputs include
\[
\text{counterfactual demand},
\qquad
\text{consumer surplus},
\qquad
\text{substitution patterns},
\]
and equilibrium outcomes.

\vspace{0.25em}

The payoff is greater extrapolation.
The cost is stronger modeling assumptions.
\end{frame}

% =========================================================
\begin{frame}{Example II: the complete workflow}

\small

\begin{tabular}{B{0.22\textwidth} L{0.68\textwidth}}
\toprule
Question
&
How would demand and welfare respond to a price policy?
\\[0.5em]

Specification
&
Utility functions, alternatives, information,
choice rule, indifference, tie-breaking.
\\[0.5em]

Estimand
&
Price coefficient?
Preference distribution?
Elasticity?
Counterfactual market share?
Welfare?
\\[0.5em]

Identification
&
What variation separates price sensitivity,
heterogeneity, indifference, and tie-breaking?
\\[0.5em]

Estimation
&
MLE?
GMM?
Simulated moments?
Bayesian methods?
\\[0.5em]

Inference
&
Standard errors?
Bootstrap?
Confidence sets for counterfactuals?
\\
\bottomrule
\end{tabular}

\vspace{0.35em}

\begin{center}
\boxed{
\text{A structural model does not escape the same five-step discipline.}
}
\end{center}

\end{frame}

% =========================================================
\section{What Is a Reduced Form?}

\begin{frame}{Why ``reduced form'' causes so much confusion}

In modern applied economics, people say things like

\begin{center}
``This is a reduced-form paper.''
\end{center}

or

\begin{center}
``We first show the reduced-form effect.''
\end{center}

But the term has an older and much more precise econometric meaning.

\vspace{0.4em}

\begin{alertblock}{Important distinction}
We should distinguish

\[
\boxed{\text{reduced form as a mathematical object}}
\]
from

\[
\boxed{\text{``reduced-form empirical work'' as a research style}.}
\]
\end{alertblock}

\end{frame}

% =========================================================
\begin{frame}{The strict econometric definition}
Consider a simultaneous structural system
\[
A_0Y_i = A_1X_i + U_i,
\]
where $Y_i$ collects the endogenous variables and $X_i$ the exogenous
variables. If $A_0$ is invertible,
\[
Y_i = A_0^{-1}A_1X_i + A_0^{-1}U_i.
\]
Define $\Pi=A_0^{-1}A_1$ and $V_i=A_0^{-1}U_i$. Then
\[
\boxed{Y_i=\Pi X_i+V_i}
\]
is the \textbf{reduced form}.

\vspace{0.25em}

Every endogenous variable is expressed as a function of
exogenous variables and disturbances.
\end{frame}

% =========================================================
\begin{frame}{Example: demand and supply}

Consider

\[
Q^d
=
\alpha_0
-
\alpha_1P
+
\alpha_2Z_d
+
u_d,
\]

\[
Q^s
=
\beta_0
+
\beta_1P
+
\beta_2Z_s
+
u_s.
\]

Equilibrium requires

\[
Q^d=Q^s=Q.
\]

The two equations are \textbf{structural}:
\[
\boxed{
\text{one describes demand behavior,
the other supply behavior.}
}
\]

Price $P$ and quantity $Q$ are jointly determined endogenous variables.

\end{frame}

% =========================================================
\begin{frame}{Solving the structural system produces the reduced form}

Set demand equal to supply:
\[
\alpha_0-\alpha_1P+\alpha_2Z_d+u_d
=
\beta_0+\beta_1P+\beta_2Z_s+u_s.
\]

Therefore,
\[
P
=
\frac{
\alpha_0-\beta_0
+\alpha_2Z_d-\beta_2Z_s
+u_d-u_s
}{
\alpha_1+\beta_1
}.
\]

Hence

\[
\boxed{
P
=
\pi_0
+
\pi_d Z_d
+
\pi_s Z_s
+
v_P.
}
\]

Substituting this expression into either structural equation gives

\[
\boxed{
Q
=
\gamma_0
+
\gamma_d Z_d
+
\gamma_s Z_s
+
v_Q.
}
\]

These equations are the reduced form.

\end{frame}

% =========================================================
\begin{frame}{Structural and reduced-form parameters are different objects}

For example,
\[
\pi_d
=
\frac{\alpha_2}{\alpha_1+\beta_1}.
\]

Thus a reduced-form coefficient generally combines several structural
parameters.

\vspace{0.35em}

Observing $\pi_d$ does not automatically tell us $\alpha_1$,
the slope of demand.

To recover structural parameters we need additional restrictions,
variation, or instruments.

\vspace{0.4em}

\begin{block}{This is the classical identification problem}
\[
\boxed{
\text{Can the structural parameters be recovered from
the reduced-form distribution?}
}
\]
\end{block}

\end{frame}

% =========================================================
\begin{frame}{The word ``reduced form'' in modern applied micro}

Modern applied researchers often use the term more broadly.

Suppose $Z_i$ is experimentally or quasi-experimentally assigned.

They may estimate

\[
Y_i
=
\alpha
+
\rho Z_i
+
\varepsilon_i
\]
and call $\hat\rho$ a \textbf{reduced-form effect}.

Why?

Because the researcher is estimating the response of an observed outcome
to plausibly exogenous variation without fully modeling the behavioral
mechanism connecting $Z$ to $Y$.

\vspace{0.35em}

\begin{alertblock}{But this is a convention, not the general mathematical definition}
OLS is not synonymous with reduced form.
\end{alertblock}

\end{frame}

% =========================================================
\begin{frame}{IV makes the terminology especially transparent}

Suppose encouragement $Z_i$ affects treatment $D_i$.

The first stage is

\[
D_i
=
\pi_0+\pi_1Z_i+v_i.
\]

The outcome reduced form is

\[
Y_i
=
\rho_0+\rho_1Z_i+e_i.
\]

The Wald estimand is

\[
\frac{\rho_1}{\pi_1}.
\]

Under the appropriate IV assumptions, this may identify a LATE.

\vspace{0.35em}

Here ``reduced form'' is close to the classical meaning:
\[
\boxed{
\text{an endogenous outcome is written directly as a function
of exogenous variation}.
}
\]

\end{frame}

% =========================================================
\begin{frame}{Four common misconceptions}

\begin{enumerate}
    \item
    \textbf{``OLS means reduced form.''}

    False. OLS is an estimation method.

    \vspace{0.2em}

    \item
    \textbf{``Logit means structural.''}

    False. Logit is a probability model.
    It becomes structural only when its objects have a behavioral
    interpretation within an economic model.

    \vspace{0.2em}

    \item
    \textbf{``Structural models must be nonlinear or complicated.''}

    False. A structural relationship can be linear.

    \vspace{0.2em}

    \item
    \textbf{``Reduced form means atheoretical.''}

    False. A credible reduced-form design requires theory about
    assignment, counterfactuals, institutions, and the estimand.
\end{enumerate}

\end{frame}

% =========================================================
\begin{frame}{So how do we tell them apart in practice?}

The misconceptions share one lesson: the econometric \emph{tool}
does not settle the classification.

\vspace{0.3em}

\begin{alertblock}{First, accept a premise}
This is not a binary classification. It is a \textbf{spectrum},
and many of the best papers sit somewhere in the middle.
\end{alertblock}

\[
\text{design-based}
\;\xleftrightarrow[\text{many good papers live in between}]{\hspace{6cm}}\;
\text{fully structural}
\]

The useful question is not ``which estimator?'' but
\[
\boxed{\text{What role does the economic model play in the paper?}}
\]

Four diagnostic questions:
\begin{enumerate}
    \item Where does the estimating equation come from?
    \item What is the target parameter?
    \item What does identification rely on?
    \item What are the estimates used for?
\end{enumerate}

\end{frame}

% =========================================================
\begin{frame}{Q1: Where does the estimating equation come from?}

\begin{columns}[T]

\column{0.48\textwidth}
\begin{block}{Reduced form}
Dictated mainly by the \textbf{identification strategy}:
the DiD double difference, the comparison around an RD cutoff.

\vspace{0.2em}

There may be a ``theoretical framework'' section for intuition.
Delete it, and the regression still runs and the result can still
be interpreted.
\end{block}

\column{0.48\textwidth}
\begin{block}{Structural}
Derived from an \textbf{explicit economic model}:
objectives, constraints, information, equilibrium conditions.

\vspace{0.2em}

The estimating equations are its first-order conditions,
likelihood, or moment conditions.
\end{block}

\end{columns}

\vspace{0.5em}

\begin{alertblock}{A useful test: take the model away}
Do the estimates still mean anything?
\begin{itemize}
    \item Yes $\Rightarrow$ the paper leans reduced form.
    \item The parameters cannot even be defined without the model
    $\Rightarrow$ structural.
\end{itemize}
\end{alertblock}

In our examples: $\bar Y_1-\bar Y_0$ survives the test;
$\alpha$ as a ``marginal utility of money'' does not.

\end{frame}

% =========================================================
\begin{frame}{Q2: What is the target parameter?}

\begin{columns}[T]

\column{0.48\textwidth}
\begin{block}{Reduced form}
Treatment effects, elasticities:
``in \emph{this} setting, when $X$ changes, how does $Y$ change?''

\vspace{0.2em}

A composite of many primitives \emph{and} the environment.
\end{block}

\column{0.48\textwidth}
\begin{block}{Structural}
\textbf{Primitives}: preferences, technology, costs,
discount rates, risk aversion, beliefs.

\vspace{0.2em}

In principle, \emph{invariant to the policy} being evaluated.
\end{block}

\end{columns}

\vspace{0.5em}

Recall the demand--supply example:
\[
\pi_d=\frac{\alpha_2}{\alpha_1+\beta_1},
\qquad
\begin{array}{l}
\alpha_1,\alpha_2:\ \text{demand primitives},\\
\beta_1:\ \text{supply primitive}.
\end{array}
\]
$\pi_d$ is a legitimate causal effect of $Z_d$ on $P$ \emph{in this market}.
A policy that changes the supply slope $\beta_1$ changes $\pi_d$,
but leaves $\alpha_1,\alpha_2$ untouched.

\vspace{0.3em}

\begin{alertblock}{Why this matters}
Reduced-form effects describe this environment;
primitives can be carried to a new one.
\end{alertblock}

\end{frame}

% =========================================================
\begin{frame}{Q3: What does identification rely on?}

\begin{columns}[T]

\column{0.48\textwidth}
\begin{block}{Reduced form}
Mainly the \textbf{research design}:
exogenous variation in the data
(randomization, discontinuities, timing of reforms).
\end{block}

\column{0.48\textwidth}
\begin{block}{Structural}
To a large extent the \textbf{model itself}:
functional forms, distributional assumptions,
equilibrium conditions, exclusion restrictions.
\end{block}

\end{columns}

\vspace{0.5em}

\textbf{Caveat.} Good structural papers \emph{also} use quasi-experimental
variation. So the question is not ``design or model?'' but
\[
\boxed{\text{How much of the identification burden does the model carry?}}
\]

\vspace{0.2em}

In our examples:
\begin{itemize}
    \item Example I: randomization does essentially all the work.
    \item Example II: with only $(D_i,X_i,p_i)$, separating
    $\alpha,\beta,\kappa,r$ leans heavily on the model,
    unless we add randomized prices, randomized defaults,
    or an explicit ``indifferent'' option.
\end{itemize}

\end{frame}

% =========================================================
\begin{frame}{Q4: What are the estimates used for?}

This is often the most visible signal: look at where the paper \emph{ends}.

\vspace{0.3em}

\begin{columns}[T]

\column{0.48\textwidth}
\begin{block}{Reduced form}
The endpoint is the \textbf{effect itself}:
report it and interpret it.
\end{block}

\column{0.48\textwidth}
\begin{block}{Structural}
The endpoint is \textbf{counterfactual simulation}:
\begin{itemize}
    \item policies never observed in the data;
    \item welfare analysis;
    \item decomposing mechanisms.
\end{itemize}
\end{block}

\end{columns}

\vspace{0.5em}

In Example II the payoff was
\[
\Pp_{\hat\theta}\bigl(D_i(p_A')=A\bigr),
\qquad
p_A'=p_A-5,
\]
a price that may never appear in the data.

\vspace{0.3em}

\begin{alertblock}{Use the four questions together}
Ask all four about any paper and you can place it on the spectrum,
whatever estimator it happens to use.
\end{alertblock}

\end{frame}

% =========================================================
\begin{frame}{Summary: the two ends of the spectrum}

\small
\setlength{\tabcolsep}{4pt}
\renewcommand{\arraystretch}{1.05}

\begin{tabular}{L{0.20\textwidth} L{0.37\textwidth} L{0.37\textwidth}}
\toprule
&
\textbf{Reduced-form / design-based}
&
\textbf{Structural}
\\
\midrule

Estimating equation
&
Dictated by the research design
&
Derived from an explicit economic model
\\[0.4em]

Target parameter
&
Treatment effects, elasticities in this setting
&
Primitives: preferences, technology, costs, beliefs
\\[0.4em]

Identification
&
Mainly exogenous variation in the data
&
Model structure + identifying variation
\\[0.4em]

End use
&
Report and interpret the effect
&
Counterfactuals, welfare, mechanisms
\\[0.4em]

Strength
&
Transparent connection between variation and estimand
&
Mechanisms and richer counterfactuals
\\[0.4em]

Main risk
&
External validity / limited counterfactual scope
&
Misspecification / stronger extrapolation assumptions
\\
\bottomrule
\end{tabular}

\vspace{0.4em}

\begin{center}
Real papers mix the two columns. The question is \emph{how much work the model does}.
\end{center}

\end{frame}

% =========================================================
\begin{frame}{Same empirical setting, different questions}

Suppose a city randomly offers a \$5 public-transit subsidy.

\textbf{Reduced-form question:} the effect
\[
\E[Y_i(1)-Y_i(0)],
\]
where treatment means receiving the subsidy offer. We might ask:
\[
\boxed{
\text{How much did ridership change because of this program?}
}
\]

\vspace{0.25em}

\textbf{Structural question:} estimate preferences over travel time, money,
comfort, and modes. Then ask:
\[
\boxed{
\begin{array}{c}
\text{What would happen under a \$2, \$10, or nonlinear subsidy,}\\
\text{or after adding a new transit mode?}
\end{array}
}
\]

Same market. Different estimands. Different assumptions.

\end{frame}

% =========================================================
\begin{frame}{Reduced form versus structural is not a contest}

It is tempting to ask:

\begin{center}
\Large
\textbf{Which approach is better?}
\end{center}

That is usually the wrong question.

\vspace{0.4em}

The correct question is:
\[
\boxed{
\text{What do we want to learn,
and what assumptions are required to learn it?}
}
\]

\vspace{0.4em}

Structural work can use randomized experiments.

Reduced-form evidence can discipline structural models.

Experimental moments can identify structural parameters.

Structural models can clarify what a reduced-form coefficient means.

\vspace{0.35em}

\begin{block}{Modern empirical work increasingly combines the two}
The boundary is methodological, not ideological.
\end{block}

\end{frame}

% =========================================================
\section{Fixed Design or Random Design?}

\begin{frame}{Two stories about the same regression}

Every course writes down the same equation:
\[
Y_i = X_i'\beta+\varepsilon_i.
\]

But different courses tell different stories about \emph{where the randomness comes from}.

\vspace{0.3em}

\begin{columns}[T]

\column{0.48\textwidth}
\begin{block}{Story 1: fixed design}
``$X$ is fixed in repeated samples.''
\[
Y_i=x_i'\beta+\varepsilon_i,
\qquad
x_i\ \text{nonrandom}.
\]
Only $\varepsilon_i$ is random.
\end{block}

\column{0.48\textwidth}
\begin{block}{Story 2: random design}
``$(Y_i,X_i)$ is an i.i.d.\ draw from a population.''
\[
(Y_i,X_i)\sim F,
\qquad
X_i\ \text{random}.
\]
Both are random.
\end{block}

\end{columns}

\vspace{0.3em}

\begin{alertblock}{And yet nearly every formula carries a $\mid X$}
\[
\E[\hat\beta\mid \bm X],
\qquad
\mathrm{Var}(\hat\beta\mid \bm X),
\qquad
\E[\varepsilon_i\mid X_i]=0.
\]
Why do we condition on something we just called fixed?
\end{alertblock}

\end{frame}

% =========================================================
\begin{frame}{Fixed design: the experimenter chooses $X$}

The fixed-design convention was inherited from experimental statistics.

In an agricultural field trial the researcher \emph{chooses} the fertilizer doses
\[
x_1,\dots,x_n,
\]
assigns them to plots, and observes yields. Here ``repeated samples'' is literal:
run the experiment again \emph{at the same $X$}.

\vspace{0.3em}

Then there is nothing to condition on, and the classical results are exact:
\[
\E[\hat\beta]=\beta,
\qquad
\mathrm{Var}(\hat\beta)=\sigma^2(\bm X'\bm X)^{-1}.
\]

This is the world in which the Gauss--Markov theorem is usually stated.

\vspace{0.35em}

\pause

\begin{alertblock}{The difficulty for economists}
We almost never choose $X$.

Nobody assigns years of schooling, firm size, or a country's institutions.
\end{alertblock}

\end{frame}

% =========================================================
\begin{frame}{Random design: what actually changes}

Under random sampling, ``another sample'' means \textbf{drawing new units},
whose regressor values would themselves be different.

So
\[
\bm X'\bm X,
\qquad
(\bm X'\bm X)^{-1},
\qquad
\hat\beta
\]
are all random, and in general
\[
\boxed{
\E\left[(\bm X'\bm X)^{-1}\right]
\;\neq\;
\left(\E[\bm X'\bm X]\right)^{-1}.
}
\]

\vspace{0.25em}

Two questions now come apart:
\[
\begin{array}{ll}
\text{Conditional:}
&
\text{how would }\hat\beta\text{ vary across samples with \emph{this} }\bm X?
\\[0.3em]
\text{Unconditional:}
&
\text{how would }\hat\beta\text{ vary across samples with a \emph{new} }\bm X?
\end{array}
\]

\begin{alertblock}{These are different reference sets}
They are two different thought experiments, and they have two different answers.
\end{alertblock}

\end{frame}

% =========================================================
\begin{frame}{Conditioning on $X$ restores the fixed-design algebra}

Write
\[
\hat\beta = \beta + (\bm X'\bm X)^{-1}\bm X'\varepsilon.
\]

Conditional on $\bm X$ this is linear in $\varepsilon$, exactly as in the fixed design.
So if
\[
\E[\varepsilon\mid \bm X]=0,
\]
then
\[
\E[\hat\beta\mid \bm X]=\beta,
\qquad
\mathrm{Var}(\hat\beta\mid \bm X)=\sigma^2(\bm X'\bm X)^{-1}
\quad\text{under }\E[\varepsilon\varepsilon'\mid \bm X]=\sigma^2 I,
\]
and by the law of iterated expectations,
\[
\E[\hat\beta]=\E\big[\E[\hat\beta\mid \bm X]\big]=\beta.
\]

\vspace{0.3em}

\begin{alertblock}{How to read the textbook sentence}
``Treat $X$ as fixed'' is an instruction to \textbf{condition on} $X$.

It is not a claim that $X$ is physically nonrandom.
\end{alertblock}

\end{frame}

% =========================================================
\begin{frame}{Is conditioning on $X$ legitimate?}

Under random sampling the joint density factors:
\[
f(y_i,x_i;\beta,\gamma)
=
\underbrace{f(y_i\mid x_i;\beta)}_{\text{contains }\beta}
\cdot
\underbrace{f(x_i;\gamma)}_{\text{does not}}.
\]

If $\beta$ and $\gamma$ are variation free, the marginal distribution of $X$ carries
no information about $\beta$: $X$ is \textbf{ancillary}. The conditionality principle
then says inference about $\beta$ should be carried out in the conditional model,
i.e.\ treating the realized $\bm X$ as the reference set.

\vspace{0.3em}

\begin{block}{When this argument fails}
\begin{itemize}
    \item $X$ is determined jointly with $Y$ (simultaneity): then $f(x)$ depends on $\beta$.
    \item Sampling depends on $Y$ (choice-based or endogenous stratified sampling).
    \item The target population differs from the sampled one, so the estimand averages
    over a \emph{different} distribution of $X$.
\end{itemize}
\end{block}

\end{frame}

% =========================================================
% [Commented out for Session 1: Examples A--D (exogeneity, LDV, variance, design-based). Uncomment to restore.]
% \begin{frame}{Example A: two exogeneity conditions}
% 
% Under a fixed design there is only one sensible condition, $\E[\varepsilon_i]=0$.
% Under a random design, two distinct conditions appear:
% \[
% \begin{array}{ll}
% \text{(i) Orthogonality:}
% &
% \E[X_i\varepsilon_i]=0,
% \\[0.3em]
% \text{(ii) Conditional mean:}
% &
% \E[\varepsilon_i\mid X_i]=0.
% \end{array}
% \]
% 
% \pause
% 
% Condition (i) holds \textbf{by construction} for the best linear predictor: it is the
% first-order condition of $\min_b\E[(Y_i-X_i'b)^2]$. It buys consistency for the
% linear projection coefficient, and nothing more.
% 
% \vspace{0.25em}
% 
% Condition (ii) is a genuine restriction. It says
% \[
% \boxed{\E[Y_i\mid X_i]=X_i'\beta,}
% \]
% i.e.\ the CEF is linear. It buys conditional unbiasedness and the interpretation of
% $\beta$ as a feature of the conditional expectation function.
% 
% \vspace{0.3em}
% 
% \begin{alertblock}{Where the interesting econometrics lives}
% Almost every identification argument in this course is an argument about (ii),
% not (i). The fixed-design story hides the distinction completely.
% \end{alertblock}
% 
% \end{frame}
% 
% % =========================================================
% \begin{frame}{Example B: strict exogeneity and the lagged dependent variable}
% 
% With i.i.d.\ sampling, $\E[\varepsilon_i\mid X_i]=0$ plus independence across units
% delivers the stronger condition
% \[
% \E[\varepsilon_i\mid \bm X]=0
% \qquad(\textbf{strict exogeneity}),
% \]
% which is what conditional unbiasedness actually requires. With dependent data it
% does not. Consider
% \[
% Y_t=\rho Y_{t-1}+\varepsilon_t,
% \qquad
% \E[\varepsilon_t\mid Y_{t-1}]=0.
% \]
% 
% Here $\varepsilon_t$ is uncorrelated with the \emph{current} regressor, but it is
% correlated with \emph{future} regressors $Y_t,Y_{t+1},\dots$, which are also rows of
% $\bm X$. Hence
% \[
% \E[\varepsilon_t\mid \bm X]\neq0,
% \qquad
% \E[\hat\rho\mid \bm X]\neq\rho,
% \qquad
% \E[\hat\rho]\neq\rho.
% \]
% 
% \begin{alertblock}{A familiar slogan, finally explained}
% ``OLS with a lagged dependent variable is biased but consistent.''
% 
% The bias is a failure of the \emph{conditioning} story, not of orthogonality.
% The same mechanism produces Nickell bias in fixed-effects panels.
% \end{alertblock}
% 
% \end{frame}
% 
% % =========================================================
% \begin{frame}{Example C: the variance formula}
% 
% \textbf{What survives.} Under $\E[\varepsilon\mid \bm X]=0$ and conditional
% homoskedasticity,
% \[
% \E\left[s^2(\bm X'\bm X)^{-1}\right]
% =
% \E\left[\mathrm{Var}(\hat\beta\mid \bm X)\right]
% =
% \mathrm{Var}(\hat\beta).
% \]
% The textbook standard error is unbiased for the conditional \emph{and} the
% unconditional variance. Conditioning by itself is not the problem.
% 
% \pause
% 
% \vspace{0.25em}
% 
% \textbf{What does not.}
% \begin{enumerate}
%     \item The shortcut $\mathrm{Var}(\hat\beta)=\sigma^2(\E[\bm X'\bm X])^{-1}$.
%     The correct object is $\E[\sigma^2(\bm X'\bm X)^{-1}]$.
% 
%     \item Conditional homoskedasticity itself. If the regression is only the best
%     \emph{linear approximation} to a nonlinear CEF, then $\mathrm{Var}(\varepsilon_i\mid X_i)$
%     varies with $X_i$ almost mechanically, and we need
%     \[
%     \mathrm{Var}(\hat\beta\mid \bm X)
%     =
%     (\bm X'\bm X)^{-1}
%     \Big(\textstyle\sum_i X_iX_i'\sigma_i^2\Big)
%     (\bm X'\bm X)^{-1}.
%     \]
% \end{enumerate}
% 
% \begin{alertblock}{Robust standard errors are a default, not a robustness check}
% Heteroskedasticity is the generic case once $X$ is random and the linear model is
% an approximation.
% \end{alertblock}
% 
% \end{frame}
% 
% % =========================================================
% \begin{frame}{Example D: a third meaning of ``fixed''}
% 
% In a randomized experiment on a given set of $n$ units, a natural model treats
% \[
% \{Y_i(1),Y_i(0),X_i\}_{i=1}^n
% \quad\text{as \textbf{fixed}},
% \]
% and the assignment vector $\bm D$ as the only random object. ``Repeated samples''
% now means \emph{re-randomization}, not resampling.
% 
% \vspace{0.3em}
% 
% So there are three different repeated experiments:
% \[
% \begin{array}{ll}
% \text{classical fixed design:}
% &
% \text{redraw }\varepsilon\text{ at the same }X,
% \\[0.25em]
% \text{random design / sampling:}
% &
% \text{redraw }(Y_i,X_i)\text{ from a population},
% \\[0.25em]
% \text{design-based:}
% &
% \text{re-randomize }\bm D\text{ on the same units}.
% \end{array}
% \]
% 
% \begin{alertblock}{They are not the same standard error}
% Neyman's randomization variance for the \emph{sample} average treatment effect is a
% different object from the sampling variance of a super-population estimand. The
% robust (EHW) variance is generally conservative for design-based uncertainty.
% \end{alertblock}
% 
% \end{frame}

% =========================================================
% [Commented out for Session 1: sample vs. population estimand. Uncomment to restore.]
% \begin{frame}{Conditioning on $X$ is part of the estimand}
% 
% Let $\tau(x)=\E[Y_i(1)-Y_i(0)\mid X_i=x]$. Two different targets:
% \[
% \boxed{
% \tau_{\text{sample}}
% =
% \frac1n\sum_{i=1}^n\tau(X_i)
% }
% \qquad\text{versus}\qquad
% \boxed{
% \tau_{\text{pop}}
% =
% \E[\tau(X_i)].
% }
% \]
% 
% The first averages over the covariate distribution we actually drew;
% the second over the population covariate distribution.
% 
% \vspace{0.25em}
% 
% Often the same estimator targets both and the point estimates coincide. But
% \[
% \mathrm{Var}\left(\hat\tau-\tau_{\text{sample}}\right)
% \;\leq\;
% \mathrm{Var}\left(\hat\tau-\tau_{\text{pop}}\right),
% \]
% because the second also pays for the fact that a new sample would have a
% different $\bm X$.
% 
% \vspace{0.3em}
% 
% \begin{alertblock}{The point of this section}
% \[
% \boxed{
% \text{``Do we condition on } X\text{?'' is a specification question,
% not a standard-error option.}
% }
% \]
% It changes what we claim to have learned, not only how precisely we learned it.
% \end{alertblock}
% 
% \end{frame}

% =========================================================
\begin{frame}{A historical note}

The fixed-design convention entered econometrics from experimental statistics,
where the researcher really does choose $X$. Economists inherited the algebra
without inheriting the experiment.

\vspace{0.3em}

For a long time ``$X$ is fixed in repeated samples'' was taught as an assumption
about the world. It is better understood as a choice of reference set:
\[
\boxed{\text{which repeated experiment are we imagining?}}
\]

\vspace{0.3em}

Modern treatments restate the linear model under random sampling, derive results
conditionally and then unconditionally, and make heteroskedasticity-robust
inference the default. The design-based literature adds a third reference set:
randomization itself.

\vspace{0.35em}

\begin{block}{Back to the workflow}
``Where does the randomness come from?'' belongs to \Spec,

long before anyone types \texttt{reg y x, robust}.
\end{block}

\end{frame}

% =========================================================
\section{Putting the Framework Together}

\begin{frame}{The same five questions apply to both traditions}

\small
\setlength{\tabcolsep}{4pt}
\renewcommand{\arraystretch}{1.1}

\begin{tabular}{L{0.18\textwidth} L{0.38\textwidth} L{0.38\textwidth}}
\toprule
&
\textbf{Treatment-effect example}
&
\textbf{Binary-choice example}
\\
\midrule

Question
&
Does vaccination change infection risk?
&
How does price affect choice and welfare?
\\[0.5em]

Specification
&
Potential outcomes, treatment, interference, horizon
&
Utility, information, choice set, tie-breaking
\\[0.5em]

Estimand
&
ATE, direct effect, spillover, policy effect
&
Preferences, elasticity, welfare, counterfactual share
\\[0.5em]

Identification
&
Randomization / conditional independence / design
&
Exclusion restrictions, price variation, functional restrictions
\\[0.5em]

Estimation
&
Difference in means, regression, weighting
&
MLE, GMM, simulated moments
\\[0.5em]

Inference
&
SEs, clustering, randomization inference
&
Asymptotic SEs, bootstrap, counterfactual uncertainty
\\
\bottomrule
\end{tabular}

\end{frame}

% =========================================================
\begin{frame}{Two different ways a project can fail}

\begin{columns}[T]

\column{0.48\textwidth}

\begin{block}{Failure 1: Wrong identification}
The estimand is meaningful, but the data do not identify it.

For example, $\ATE$ is the right question, but treatment selection
is endogenous.

\vspace{0.35em}

Problem:
\[
\boxed{\text{counterfactual is not identified}}
\]

\end{block}

\column{0.48\textwidth}

\begin{block}{Failure 2: Wrong specification}
The parameter is perfectly identified, but it is not the object
corresponding to the substantive question.

For example, estimate $\E[Y(1)-Y(0)]$ when the policy operates
through network spillovers.

\vspace{0.35em}

Problem:
\[
\boxed{\text{we precisely answer the wrong question}}
\]

\end{block}

\end{columns}

\end{frame}

% =========================================================
\begin{frame}{A useful hierarchy of econometric questions}

When reading a paper, do not begin with:

\begin{center}
``What regression did they run?''
\end{center}

Instead ask:

\vspace{0.25em}

\begin{enumerate}
    \item What is the substantive question?
    \item What mathematical object represents the answer?
    \item What assumptions make that object meaningful?
    \item What assumptions connect it to observed data?
    \item What variation actually identifies it?
    \item What estimator implements that identification argument?
    \item What uncertainty does the reported standard error capture?
    \item What important uncertainty does it \emph{not} capture?
\end{enumerate}

\end{frame}

% =========================================================
\begin{frame}{When you see a regression, ask what role it plays}
Consider
\[
Y_i = \alpha+\beta D_i+X_i'\gamma+\varepsilon_i.
\]
The same equation can play very different roles. It might be
\[
\begin{array}{c}
\boxed{\text{a descriptive conditional expectation}},\\[0.3em]
\boxed{\text{an estimator implementing a randomized experiment}},\\[0.3em]
\boxed{\text{a selection-on-observables identification strategy}},\\[0.3em]
\boxed{\text{one equation inside a structural model}}.
\end{array}
\]

\begin{alertblock}{The syntax of the regression does not reveal its scientific meaning}
The meaning comes from the specification and identification argument.
\end{alertblock}
\end{frame}

% =========================================================
\begin{frame}{The regression coefficient is not the starting point}

A common workflow is:
\[
\text{write regression}
\rightarrow
\text{estimate }\beta
\rightarrow
\text{interpret }\beta.
\]

A better workflow is:
\[
\boxed{
\text{Question}
\rightarrow
\text{Estimand}
\rightarrow
\text{Identification}
}
\]
and only then

\[
\boxed{
\text{choose a statistical procedure that estimates the identified object}.
}
\]

\vspace{0.4em}

Modern developments in

\[
\text{IV},
\quad
\text{DiD},
\quad
\text{RDD},
\quad
\text{treatment effects},
\quad
\text{structural estimation}
\]
are much easier to understand from this perspective.

\end{frame}

% =========================================================
\begin{frame}{A warning about ``assumptions''}

Econometric assumptions play different roles.

\vspace{0.3em}

\textbf{Specification assumptions:}

\[
Y_i(d),
\qquad
U_{ij}=X_{ij}'\beta+\varepsilon_{ij},
\qquad
\text{choice set},
\qquad
\text{equilibrium concept}.
\]

These determine what world we are modeling.

\vspace{0.3em}

\textbf{Identification assumptions:} $D\perp(Y(1),Y(0))$,
exclusion restrictions, rank conditions, monotonicity,
support conditions, etc.

These determine what can be learned from observables.

\vspace{0.3em}

\textbf{Inference assumptions:}

sampling structure, clustering, dependence,
regularity conditions.

\vspace{0.35em}

\begin{alertblock}{Do not put all assumptions into one undifferentiated list}
They answer different logical questions.
\end{alertblock}

\end{frame}

% =========================================================
\begin{frame}{What does ``credible'' mean?}
A credible empirical result requires alignment across all stages:
\[
\boxed{
\begin{array}{ccccc}
\text{Question}
&
\longleftrightarrow
&
\text{Estimand}
&
\longleftrightarrow
&
\text{Specification}
\\
&&\downarrow&&\\
&&\text{Identification}&&\\
&&\downarrow&&\\
&&\text{Estimation}&&\\
&&\downarrow&&\\
&&\text{Inference}&&
\end{array}
}
\]
A sophisticated estimator cannot compensate for
\[
\text{a poorly defined question},
\quad
\text{an inappropriate model},
\]
or
\[
\text{an unconvincing counterfactual}.
\]
\end{frame}

% =========================================================
\begin{frame}{The main methodological habit}

Before touching the data, try to complete the following sentence:

\vspace{0.4em}

\begin{center}
\Large
\textit{``I want to learn \underline{\hspace{3cm}},
\\[0.4em]
and I can learn it from these data because
\underline{\hspace{3cm}}.''}
\end{center}

\vspace{0.5em}

The first blank is primarily about

\[
\boxed{\text{Question + Specification + Estimand}.}
\]

The second blank is primarily about

\[
\boxed{\text{Identification}.}
\]

\vspace{0.35em}

Only afterward should we ask:
\[
\boxed{\text{How should I estimate it?}}
\]

\end{frame}

% =========================================================
\begin{frame}{A roadmap for the rest of the course}

We will repeatedly use the same framework.

\[
\boxed{
\text{Question}
\rightarrow
\text{Specification \& Estimand}
\rightarrow
\text{Identification}
\rightarrow
\text{Estimation}
\rightarrow
\text{Inference}
}
\]

Examples will include:
\[
\begin{array}{c}
\text{Regression and selection on observables}\\
\text{Instrumental variables and LATE}\\
\text{Difference-in-differences}\\
\text{Regression discontinuity}\\
\text{Panel data}\\
\text{Treatment-effect heterogeneity}\\
\text{Structural models and policy counterfactuals}
\end{array}
\]

The methods will change.

The sequence of questions should not.

\end{frame}

% =========================================================
\begin{frame}{Takeaways}

\begin{enumerate}
    \item
    Econometrics begins with a \textbf{question}, not an estimator.

    \vspace{0.25em}

    \item
    An \textbf{estimand} is the precise mathematical object that answers
    the question.

    \vspace{0.25em}

    \item
    \textbf{Specification} determines what world we are modeling.

    \vspace{0.25em}

    \item
    \textbf{Identification} determines what can be learned from
    the population distribution of observables.

    \vspace{0.25em}

    \item
    \textbf{Estimation} and \textbf{inference} come afterward.

    \vspace{0.25em}

    \item
    Reduced-form and structural methods are not rival religions but
    two ends of a spectrum: classify a paper by the role its model
    plays, not by its estimator.
\end{enumerate}

\end{frame}

% =========================================================
\begin{frame}{One sentence to remember}

\begin{center}

\vspace{0.6em}

\Huge
\textbf{
Do not ask
\\[0.3em]
``Which estimator should I use?''
\\[0.8em]
until you know
\\[0.3em]
``What exactly am I trying to learn?''
}

\end{center}

\end{frame}

% =========================================================
\begin{frame}{Selected references}

\scriptsize

Angrist, J. D., and J.-S. Pischke (2009).
\textit{Mostly Harmless Econometrics: An Empiricist's Companion}.
Princeton University Press.

\vspace{0.25em}

Heckman, J. J. (2010).
``Building Bridges between Structural and Program Evaluation
Approaches to Evaluating Policy.''
\textit{Journal of Economic Literature}, 48(2), 356--398.

\vspace{0.25em}

Heckman, J. J., and E. Vytlacil (2005).
``Structural Equations, Treatment Effects and Econometric Policy
Evaluation.''
\textit{Econometrica}, 73(3), 669--738.

\vspace{0.25em}

Low, H., and C. Meghir (2017).
``The Use of Structural Models in Econometrics.''
\textit{Journal of Economic Perspectives}, 31(2), 33--58.

\vspace{0.25em}

Hudgens, M. G., and M. E. Halloran (2008).
``Toward Causal Inference with Interference.''
\textit{Journal of the American Statistical Association},
103(482), 832--842.

\vspace{0.25em}

White, H. (1980).
``A Heteroskedasticity-Consistent Covariance Matrix Estimator and a Direct
Test for Heteroskedasticity.''
\textit{Econometrica}, 48(4), 817--838.

\vspace{0.25em}

Abadie, A., S. Athey, G. W. Imbens, and J. M. Wooldridge (2020).
``Sampling-Based versus Design-Based Uncertainty in Regression Analysis.''
\textit{Econometrica}, 88(1), 265--296.

\end{frame}

% =========================================================
\begin{frame}

\begin{center}

\vspace{1em}

\LARGE
\textbf{
We choose to learn econometrics,
\\[0.5em]
not because it is easy,
\\[0.3em]
but because it is hard.
}

\vspace{1.5em}

\normalsize
\textcolor{DarkGray}{--- after John F. Kennedy, Rice University, 1962}

\end{center}

\end{frame}

\end{document}